\documentclass[11pt,a4paper]{article}
\usepackage[T1]{fontenc}
\usepackage[utf8]{inputenc}
\usepackage[main=italian,provide=*]{babel}
\usepackage{amsmath,amssymb}
\usepackage{geometry}
\geometry{margin=2.3cm,top=2.5cm,bottom=2.7cm}
\usepackage{booktabs}
\usepackage[table]{xcolor}
\usepackage{tcolorbox}
\tcbuselibrary{skins,breakable}
\usepackage{titlesec}
\usepackage{enumitem}
\usepackage{seqsplit}
\usepackage{needspace}
\usepackage{fancyhdr}
\usepackage{hyperref}
\hypersetup{colorlinks=true,linkcolor=Sepia,citecolor=Sepia,urlcolor=Sepia}

% --- palette ---
\definecolor{inkblue}{HTML}{1B3A5C}
\definecolor{lightblue}{HTML}{EAF1F8}
\definecolor{accent}{HTML}{B0562A}
\definecolor{softgray}{HTML}{6B6B6B}
\definecolor{okgreen}{HTML}{2E6B3E}
\definecolor{okgreenbg}{HTML}{EAF5EC}
\definecolor{warnamber}{HTML}{9C6B12}
\definecolor{warnbg}{HTML}{FBF1E0}
\definecolor{notebar}{HTML}{9C6B12}

% --- titoli di sezione ---
\titleformat{\section}
  {\normalfont\Large\bfseries\color{inkblue}}
  {\thesection}{0.6em}{}[{\color{inkblue!40}\titlerule}]
\titlespacing{\section}{0pt}{0.8em}{0.4em}
\titleformat{\subsection}
  {\normalfont\large\bfseries\color{accent}}
  {}{0em}{}
\titlespacing{\subsection}{0pt}{0.5em}{0.2em}

\pagestyle{fancy}
\fancyhf{}
\renewcommand{\headrulewidth}{0.4pt}
\fancyhead[L]{\small\color{softgray}Note su \texttt{confronto\_ansatz\_entangler.ipynb}}
\fancyhead[R]{\small\color{softgray}\thepage}
\fancyfoot[C]{\small\color{softgray}Tesi triennale, Samuele Schiavi --- Relatore: Prof.\ Alessandro Chiesa}

% --- box riutilizzabili ---
\newtcolorbox{keyeq}[1][]{
  colback=lightblue, colframe=inkblue!70, boxrule=0.6pt,
  arc=2pt, left=8pt, right=8pt, top=4pt, bottom=4pt, #1
}
\newtcolorbox{panelbox}[1]{
  colback=white, colframe=softgray!50, boxrule=0.5pt,
  arc=2pt, left=7pt, right=7pt, top=3pt, bottom=3pt,
  fonttitle=\bfseries\color{inkblue}, title={#1}
}
\newtcolorbox{okbox}[1]{
  colback=okgreenbg, colframe=okgreen!60, boxrule=0.6pt,
  arc=2pt, left=7pt, right=7pt, top=3pt, bottom=3pt,
  fonttitle=\bfseries\color{okgreen}, title={#1}
}
\newtcolorbox{warnbox}[1]{
  colback=warnbg, colframe=warnamber!70, boxrule=0.6pt,
  arc=2pt, left=7pt, right=7pt, top=3pt, bottom=3pt,
  fonttitle=\bfseries\color{warnamber}, title={#1}
}
\newtcolorbox{editnote}{
  colback=white, colframe=white, boxrule=0pt,
  borderline west={2.2pt}{0pt}{notebar},
  arc=0pt, left=10pt, right=4pt, top=2pt, bottom=2pt,
  fontupper=\itshape\small\color{softgray!90!black}
}

\title{\vspace{-1.5em}\color{inkblue}Note sul notebook\\[0.1em]\texttt{confronto\_ansatz\_entangler.ipynb}}
\author{Note di lavoro --- Tesi triennale, Samuele Schiavi\\ Relatore: Prof.\ Alessandro Chiesa}
\date{}

\begin{document}
\sloppy
\maketitle
\thispagestyle{fancy}
\vspace{-2em}

Il notebook confronta \textbf{quattro ansatz variazionali con entangler
diversi} (A--D), tutti costruiti con sole rotazioni $R_y$ (quindi in
campo reale), più \textbf{due famiglie PMA} a rottura di simmetria. Ogni
ansatz è ottimizzato via VQE su simulatore statevector (sistema chiuso,
esatto) con la metodologia adottata: \textbf{multistart con
reinizializzazione casuale ad ogni ciclo esterno}. Poggia sullo stesso
benchmark esatto (\texttt{eigh}) usato dagli altri moduli del dimero.

\begin{center}
\renewcommand{\arraystretch}{1.3}
\small
\begin{tabular}{@{}lp{3.0cm}p{2.6cm}cp{2.6cm}@{}}
\toprule
\rowcolor{inkblue!10}
 & \textbf{Circuito} & \textbf{Entangler} & \textbf{N.\ par.} & \textbf{Note}\\
\midrule
\textbf{A} & $R_y$--CNOT--$R_y$ & CNOT (non param.) & 4 & ridondanza 1\\
\textbf{B} & $R_y$--CZ--$R_y$ & CZ (non param.) & 4 & ridondanza 1, $\equiv$ A\\
\textbf{C} & $R_y$--$CR_y(\varphi)$--$R_y$ & $CR_y$ (param.) & 5 & ridondanza 2\\
\textbf{D} & $R_y$--RBS($\varphi$)--$R_y$ & RBS (param.) & 5 & ridondanza 2\\
\textbf{PMA-1q} & $X$--RBS--$[R_y(q_1)$--RBS$]^k$ & RBS + $R_y$ su un solo qubit & 1,2,4,6 & $\mathcal F{=}1$ da 4\\
\textbf{PMA-2q} & $X$--RBS--$[R_y(q_0),R_y(q_1)]^k$ & RBS + $R_y$ indip.\ su entrambi & 1,3,5,7 & \textbf{$\bigstar$\ consigliata}, $\mathcal F{=}1$ da 3\\
\bottomrule
\end{tabular}
\end{center}

La \textbf{ridondanza} di A--D è $n_\text{par}-3$: il ground state reale
di 2 qubit vive in $S^3$, cioè ha 3 gradi di libertà fisici (4 ampiezze
reali $-1$ normalizzazione $-1$ fase globale). Ogni parametro oltre i 3
è gauge. Per le famiglie PMA il conteggio è diverso: si parte da 1 solo
parametro (il blocco RBS di base, $M$-conservante) e si aggiungono
parametri di \emph{rottura di simmetria}, uno alla volta (PMA-1q) o a
coppie indipendenti (PMA-2q).

\section{Un solo blocco $M$-conservante in tutto il notebook: il blocco RBS}

Sia l'ansatz D sia le famiglie PMA usano la \textbf{stessa, unica
implementazione} del blocco $M$-conservante --- niente decomposizioni
diverse per la stessa fisica, per evitare l'ambiguità di avere due gate
"equivalenti ma scritti diversamente" a spasso nel notebook.

Il blocco realizza la rotazione di Givens reale nel settore
$\{|01\rangle,|10\rangle\}$:
\begin{keyeq}
\centering
$G(\varphi)=\begin{pmatrix}1&&&\\&\cos\varphi&-\sin\varphi&\\&\sin\varphi&\cos\varphi&\\&&&1\end{pmatrix}$,
\quad con soli $R_y$, Hadamard e CZ:\\[4pt]
$\mathrm{RBS}(\varphi) = (H\otimes H)\,\mathrm{CZ}\,\big(R_y(\varphi)_{q_0}\otimes R_y(-\varphi)_{q_1}\big)\,\mathrm{CZ}\,(H\otimes H)$
\end{keyeq}
nessuna decomposizione CX-$CR_y$-CX ad hoc. Self-test nel notebook:
errore rispetto a $G(\varphi)$ target $\sim10^{-16}$, per diversi valori
di $\varphi$.

Da solo, il blocco RBS conserva $M$ (è il cuore del PMA); circondato da
strati $R_y$ liberi, come nell'ansatz D, esplora invece stati generici
--- è il \textbf{ponte HA$\leftrightarrow$PMA}: la stessa unità fisica
che, a seconda di cosa le sta intorno, si comporta da ansatz simmetrico
o da ansatz generico.

\section{Perché si può restare in campo reale (anche con $D\neq0$)}

L'Hamiltoniana
$$H = J\,(X_1X_2+Y_1Y_2+Z_1Z_2) + b\,(Z_1+Z_2) + D\,(X_1Z_2 - Z_1X_2)$$
è una \textbf{matrice reale simmetrica}: tutti i termini, incluso
$Y\otimes Y$ (i due fattori $i$ si elidono) e il DM in forma $XZ-ZX$,
hanno entrate reali. Per il teorema spettrale, \textbf{il ground state
può sempre essere scelto reale} --- cercarlo tra i soli stati reali
(ansatz a sole $R_y$) è una restrizione \emph{gratuita}, non perde
generalità.

Il notebook lo verifica in una cella: a $D=0.2$ la parte immaginaria di
$H$ è esattamente $0$ e il GS risulta reale.

\begin{editnote}
\textbf{Nota di rigore --- quale DM.} Il termine usato qui,
$D(X_1Z_2-Z_1X_2)$, è quello del notebook di partenza del progetto. Va
distinto dalla DM "canonica" $D(X_1Y_2-Y_1X_2)$:
$XZ-ZX$ è \textbf{reale}, ma \textbf{rompe la conservazione di $M$}
($\|[H,S_z]\|\neq0$, verificato in cella); $XY-YX$ è
\textbf{immaginaria}, renderebbe $H$ complesso e il GS complesso. La
distinzione è cruciale per l'interpretazione: ciò che il DM di questo
notebook rompe è la \textbf{simmetria $M$}, non la realtà dello stato.
È esattamente questo a mettere in crisi il PMA-1q, come si vede dai
risultati.
\end{editnote}

\section{Le due famiglie PMA: come rompere $M$ in modo reale}

Il PMA di base ($X\cdot\mathrm{RBS}(\theta)$, 1 parametro) resta nel
settore $M=0$ per costruzione, e questo lo fa fallire in \textbf{due
situazioni distinte}: già a $D=0$ per $b/J>2$ (il ground state è
$|11\rangle$, settore $M=-1$, semplicemente irraggiungibile con 1
parametro $M$-conservante --- non è un effetto del DM, è l'assenza di
un meccanismo per uscire da $M=0$), e a $D\neq0$ vicino a $b/J=2$ (dove
il DM mescola i settori). Per recuperare in entrambi i casi bisogna
popolare i settori $M=\pm1$ con rotazioni $R_y$ \textbf{reali}
aggiuntive. Due modi funzionano (analisi completa, con derivazioni e
casi che \emph{falliscono}, in \texttt{analisi\_rotazioni\_PMA.ipynb}):

\begin{itemize}[leftmargin=1.3em,itemsep=0.25em]
\item \textbf{PMA-1q}: $R_y$ su \textbf{un solo} qubit ($q_1$) alternato
a blocchi RBS. 1 parametro per blocco $\to$ sequenza \textbf{1, 2, 4,
6}. Funziona perché la rotazione locale agisce su uno stato \emph{già
entangled} dal blocco RBS precedente: basta un solo grado di libertà
locale per redistribuire ampiezza verso $M=\pm1$. Raggiunge
$\mathcal F=1$ da \textbf{4} parametri.
\item \textbf{PMA-2q} $\bigstar$\ \textbf{(scelta consigliata)}: $R_y$
\textbf{indipendenti} su entrambi i qubit ($\alpha\neq\beta$). 2
parametri per coppia $\to$ sequenza \textbf{1, 3, 5, 7}. Raggiunge
$\mathcal F=1$ già da \textbf{3} parametri --- il minimo teorico ($S^3$,
zero ridondanza).
\end{itemize}

\textbf{Perché non basta un parametro condiviso.} Ruotare entrambi i
qubit con lo \textbf{stesso} parametro (stesso segno o segno opposto)
non funziona: i coefficienti generati nei settori $M=\pm1$ risultano
entrambi proporzionali a un'unica combinazione ($a+b$ o $a-b$) --- un
solo grado di libertà reale, non due. Serve un parametro
\emph{indipendente} per settore. Derivazione completa in
\texttt{analisi\_rotazioni\_PMA.ipynb}.

\textbf{Scelta canonica per il seguito.} Si tengono entrambe le famiglie
nel confronto (il contrasto è istruttivo), ma \textbf{PMA-2q$\cdot$3} è
il riferimento che verrà usato nei prossimi notebook ($N=3$): a parità
di fisica, raggiunge il risultato con meno parametri totali. La
variabile \texttt{PMA\_RECOMMENDED} nel codice la marca esplicitamente
($\bigstar$\ nei titoli dei grafici).

\needspace{10\baselineskip}
\section{Multistart}

Il paesaggio
$E(\boldsymbol\theta)=\langle\psi(\boldsymbol\theta)|H|\psi(\boldsymbol\theta)\rangle$
è \textbf{non-convesso}: un ottimizzatore locale (COBYLA) converge al
minimo del bacino in cui cade il punto iniziale. Il multistart campiona
$R$ punti di partenza casuali e tiene la corsa a energia minima:
\begin{keyeq}
\centering
$\boldsymbol\theta^\star = \displaystyle\arg\min_{k=1,\dots,R}\; E\big(\boldsymbol\theta^{*(k)}\big)$
\end{keyeq}

Due precisazioni operative:
\begin{itemize}[leftmargin=1.3em,itemsep=0.25em]
\item La reinizializzazione è nel ciclo \textbf{esterno} (una per corsa
completa), \textbf{non} ad ogni iterazione interna di COBYLA: azzerare
$\boldsymbol\theta$ ad ogni passo distruggerebbe la discesa. Crippa et
al.\ usano 5 restart (fino a 10 nei casi difficili); qui $R=6$.
\item La funzione registra anche la \textbf{dispersione} delle $R$
energie finali (quanto è accidentato il paesaggio) e il costo
(\texttt{nfev} della corsa migliore). Così il notebook non solo trova
il ground state, ma \textbf{mostra perché} il multistart serve.
\end{itemize}

\begin{okbox}{Distinzione da tenere ferma}
I restart e la scelta dell'ansatz risolvono problemi \emph{ortogonali}:
l'\textbf{architettura} decide \emph{cosa} è raggiungibile (l'immagine
dell'ansatz), i \textbf{restart} aiutano a \emph{trovarlo} dentro un
paesaggio con minimi locali. Se il ground state non è raggiungibile,
nessun numero di restart lo recupera. È esattamente il caso di
PMA-1q$\cdot$1 a $D\neq0$: nessun restart lo salva, perché il settore
$M=0$ non contiene il vero ground state.
\end{okbox}

\section{Fidelity robusta alla degenerazione}

A $D=0$ e $b/J=2$ esatto, singoletto ($E=-3J$) e $|11\rangle$
($E=J-2b=-3J$) sono \textbf{degeneri}: \texttt{eigh} restituisce una
combinazione arbitraria del sottospazio fondamentale, e la fidelity
verso \emph{un} singolo autovettore è mal definita. Il notebook usa la
fidelity verso l'\textbf{intero sottospazio fondamentale} $\mathcal G$:
\begin{keyeq}
\centering
$\mathcal F = \langle\psi|\,P_{\mathcal G}\,|\psi\rangle = \displaystyle\sum_{i\in\mathcal G}|\langle v_i|\psi\rangle|^2$
\end{keyeq}
con $P_{\mathcal G}$ proiettore sugli autostati a energia $E_0$ (entro
tolleranza). Fuori dalla degenerazione si riduce alla fidelity
ordinaria.

\section{Grafici interattivi (\S 6.2, 6.3, 6.4)}

Con 11 ansatz nel confronto, un grafico statico con tutte le curve è
illeggibile. Le figure di queste tre sezioni sono \textbf{interattive}:
una griglia di checkbox (una per ansatz) più due pulsanti "Seleziona
tutti"/"Deseleziona tutti", collegata al grafico tramite
\texttt{ipywidgets.interactive\_output} --- selezionando/deselezionando
le curve, il grafico si ridisegna automaticamente, senza rieseguire la
cella. Se \texttt{ipywidgets} non è installato, la cella lo rileva
(\texttt{HAS\_WIDGETS}) e mostra automaticamente un grafico statico con
tutte le curve, senza errori.

\S 6.1 (fidelity, un pannello per ansatz) resta invece una griglia
statica: con un pannello per ansatz non c'è affollamento da diradare.

\needspace{10\baselineskip}
\section{Il risultato centrale}

\begin{center}
\renewcommand{\arraystretch}{1.35}
\small
\begin{tabular}{@{}p{2.6cm}p{2.9cm}p{3.6cm}p{3.6cm}@{}}
\toprule
\rowcolor{inkblue!10}
\textbf{Configurazione} & \textbf{Esito A--D} & \textbf{Esito PMA-1q$\cdot$1} & \textbf{Esito PMA-2q$\cdot$3}\\
\midrule
$D=0$, $b/J<2$ & $\mathcal F=1$, $\Delta E\sim10^{-9}$ & $\mathcal F=1$ (settore $M{=}0$, GS=singoletto) & $\mathcal F=1$\\
$D=0$, $b/J=2$ & energia esatta; serve fidelity-sottospazio & $\mathcal F=1$ (settore $M{=}0$) & $\mathcal F=1$\\
$D=0$, $b/J>2$ & $\mathcal F=1$ & $\mathcal F=0$ (GS=$|11\rangle$, $M{=}-1$, irraggiungibile) & $\mathcal F=1$\\
$D=0.2$--$1.0$, ogni $b/J$ & $\mathcal F=1$ ovunque & $\mathcal F=0$ per $b/J{>}2$; peggio vicino a $b/J{=}2$ & $\mathcal F=1$ ovunque\\
\bottomrule
\end{tabular}
\end{center}

\subsection*{Cinque letture}

\begin{enumerate}[leftmargin=1.4em,itemsep=0.4em]
\item \textbf{A $\equiv$ B.} Le curve di A ($R_y$-CNOT) e B ($R_y$-CZ)
coincidono: per $N=2$ CNOT e CZ, avvolti da strati $R_y$, sono
localmente equivalenti
($\mathrm{CNOT}=(I\otimes H)\,\mathrm{CZ}\,(I\otimes H)$, con le
Hadamard assorbite dagli $R_y$). La scelta dell'entangler è qui
\textbf{cosmetica} --- conta solo per la transpilazione su hardware
reale.

\item \textbf{Gli ansatz reali bastano anche a $D\neq0$.} A $D=0.2$
tutti e quattro raggiungono $\mathcal F=1$: il GS resta reale, quindi
le sole $R_y$ sono sufficienti.

\item \textbf{PMA-1q$\cdot$1 fallisce per assenza di simmetrie di
rottura, non solo per il DM.} Già a $D=0$ crolla a fidelity $0$ per
$b/J>2$: parte sempre dal riferimento $|01\rangle$ ($M=0$) e il blocco
RBS da solo non può mai raggiungere $|11\rangle$ ($M=-1$), il vero
ground state sopra soglia --- è un limite dell'assenza di branching
sullo stato iniziale, non del DM. A $D=0.2$ il quadro si allarga con un
secondo tipo di fallimento, concentrato attorno a $b/J=2$, dove il DM
mescola i settori. Il contrasto con gli ansatz reali generici --- che
\emph{non} hanno nessuno di questi due vincoli e restano a
$\mathcal F=1$ ovunque --- \textbf{isola} il messaggio: il limite del
PMA-1q$\cdot$1 è strutturale (confinamento a $M=0$), non la mancanza di
fasi complesse.

\item \textbf{La rottura di simmetria reale funziona, e in due modi
diversi.} PMA-1q e PMA-2q dimostrano che bastano rotazioni $R_y$
aggiuntive --- non serve uscire dal campo reale --- per recuperare
$\mathcal F=1$ anche a $D\neq0$. PMA-2q$\cdot$3 ci arriva con un
parametro in meno perché fornisce \textbf{due gradi di libertà
indipendenti in un colpo solo} (uno per settore $M=\pm1$), mentre
PMA-1q ne fornisce uno alla volta.

\item \textbf{Espressività vs costo.} C e D (5 param., ridondanza 2)
non migliorano la fidelity rispetto ad A/B (già a 1), ma hanno un
paesaggio più accidentato: dispersione multistart maggiore, più
valutazioni di funzione. Conferma operativa del principio
\textbf{"parametri adeguati, non massimi"}: oltre i 3 gradi di libertà
fisici si aggiunge solo gauge, che l'ottimizzatore deve comunque
navigare. Lo stesso principio, applicato alle famiglie PMA, è
dimostrato esplicitamente in \S 6.6: impilare blocchi RBS aggiuntivi
(tutti $M$-conservanti) non cambia la fidelity --- restano nel settore
$M=0$ qualunque sia $K$.
\end{enumerate}

\needspace{14\baselineskip}
\section{La lezione}

Questo notebook completa e rafforza per contrasto il risultato di
\texttt{vqe\_dimer\_spiegato.ipynb}. Là si vedeva il PMA degradare a
$D\neq0$; qui si \textbf{dimostra la causa}: mettendo accanto quattro
ansatz reali generici che non degradano affatto, si esclude la
complessità dello stato come spiegazione e si isola la \textbf{rottura
di simmetria $M$} come unico responsabile --- e si mostra \textbf{come
ripararla} in modo economico, con due famiglie PMA estese che restano
nel campo reale.

Emergono tre principi di progetto degli ansatz:
\begin{itemize}[leftmargin=1.3em,itemsep=0.25em]
\item \textbf{Raggiungibilità e minimi locali sono problemi separati}:
il primo si risolve con l'architettura, il secondo con il multistart.
Confonderli porta a "aggiungere restart" quando il problema è l'ansatz,
o "cambiare ansatz" quando basta un restart in più.
\item \textbf{La restrizione al campo reale è gratuita} finché $H$ è
reale (anche con il DM $XZ-ZX$): un ansatz a sole $R_y$ non perde
nulla, e risparmia parametri rispetto a uno con $R_z$.
\item \textbf{Non tutti i parametri aggiuntivi sono uguali}: impilare
parametri $M$-conservanti non aiuta mai (\S 6.6); servono parametri che
aprano \emph{nuovi} gradi di libertà verso i settori mancanti. È la
stessa distinzione che decide tra PMA-1q e PMA-2q, e che renderà
cruciale la scelta dell'entangler per il triangolo frustrato.
\end{itemize}

Sarà rilevante per il \textbf{triangolo frustrato} ($N=3$), dove la
scelta dell'entangler smette di essere cosmetica: sfruttare (o rompere)
la simmetria di rotazione $C_3$ della topologia cambia l'immagine
raggiungibile dell'ansatz, e la distinzione HA$\leftrightarrow$PMA ---
di cui il blocco RBS è qui il ponte esplicito --- diventa strutturale.
PMA-2q$\cdot$3 è il punto di partenza naturale per quella fase.

\end{document}
